\problemname{KATT och råtta}


\section*{Interaktion}



\textbf{Se till att flusha outputen efter varje query}, annars kan du få \textit{Time Limit Exceeded}.
I C++ kan detta göras med exempelvis \texttt{cout << flush;}
eller \texttt{fflush(stdout);},
i Python med \texttt{stdout.flush()}
och i Java med \texttt{System.out.flush();}.

\section*{Poängsättning}
Din lösning kommer att testas på en mängd testfallsgrupper.
För att få poäng för en grupp så måste du klara alla testfall i gruppen.

\noindent
\begin{tabular}{| l | l | l |}
  \hline
  Grupp & Poängvärde & Gränser \\ \hline \hline
  $1$   & $5$        & $N = 2$ \\ \hline
  $2$   & $10$        & $N \leq 10$ \\ \hline
  $3$   & $10$        & $N \leq 250$ \\ \hline
  $4$   & $16$        & Varje rum som inte ligger längst ner i källaren har två rum under sig. \\ \hline
  $5$   & $13$        & $N \leq 500$ \\ \hline
  $6$   & $21$        & $N \leq 950$ \\ \hline
  $7$   & $25$        & $N \leq 1000$ \\ \hline
\end{tabular}
